% =====================================================================
% Section III: Methodology
% A-FONF: Adaptive Fractional-Order Notch Filtering
% =====================================================================
\section{Methodology}
\label{sec:method}

\subsection{Measurement Model and Notation}
\label{sec:model}

Let $f\in\mathbb{R}_{+}^{N}$ denote the discretised attenuation image on an
$n\times n$ grid, $N=n^{2}$, and let $A\in\mathbb{R}_{+}^{M\times N}$ be the
system matrix of the acquisition, with $M=n_{d}n_{\theta}$ for $n_{d}$
detector bins and $n_{\theta}$ views. Photon-limited measurements are modelled
as independent Poisson variables
\begin{equation}
  y_{i}\sim\mathrm{Poisson}\!\left(\bar{y}_{i}\right),\qquad
  \bar{y}_{i}=\frac{(Af)_{i}}{\kappa},\qquad
  \kappa=\frac{\max_{j}(Af)_{j}}{I_{0}},
  \label{eq:poisson}
\end{equation}
where $I_{0}$ is the incident count level of the most attenuating ray and
$\kappa$ converts counts to line integrals, $g=\kappa y$. Lower $I_{0}$
therefore corresponds to a lower dose. All quantities in
\eqref{eq:poisson} are known at reconstruction time except $f$: the
calibration constant $\kappa$ follows from the scanner gain, and the counts
$y$ are the raw measurements.

Frequencies are expressed in radians per pixel, $\omega=(\omega_{x},
\omega_{y})$ with $|\omega|\le\omega_{N}\sqrt{2}$ and Nyquist radius
$\omega_{N}=\pi$. We write $\mathcal{F}$ for the two-dimensional discrete
Fourier transform, $\odot$ for the Hadamard product and $\mathbf{1}$ for the
all-ones vector.

\subsection{Fractional-Order Notch Filter}
\label{sec:filter}

The regulariser acts in the frequency domain through a real, zero-phase
transfer function that combines a fractional-order low-pass envelope with
$K$ narrowband stop bands,
\begin{equation}
  H_{\alpha,\Omega}(\omega)=
  \underbrace{\exp\!\left(-\frac{\alpha\,|\omega|^{2}}{4\omega_{N}^{2}}\right)}_{\text{envelope}}
  \prod_{m=1}^{K}
  \underbrace{\left[1-\exp\!\left(-\frac{\|\omega-\omega_{m}\|^{2}}{2\sigma_{m}^{2}}\right)\right]^{\alpha/2}}_{\text{notch }m},
  \label{eq:filter}
\end{equation}
where $\alpha>0$ is the fractional order,
$\Omega=\{(\omega_{m},\sigma_{m})\}_{m=1}^{K}$ collects the notch centres and
widths, and each factor is applied at $\omega_{m}$ and at $-\omega_{m}$ so
that $H_{\alpha,\Omega}(-\omega)=H_{\alpha,\Omega}(\omega)$. The associated
spectral operator \eqref{eq:operator} is
\begin{equation}
  \mathcal{S}_{\alpha,\Omega}\,x=\mathcal{F}^{-1}\!\left\{H_{\alpha,\Omega}\odot\mathcal{F}\{x\}\right\}.
  \label{eq:operator}
\end{equation}
Two design choices matter for what follows. First, $\alpha$ scales the
exponent of $|\omega|^{2}/\omega_{N}^{2}$, so it is defined relative to the
sampling grid and is invariant to image resolution when the sampling density
scales with $n$. Second, the notch depth carries the same order $\alpha$, so
a single parameter controls both the broadband envelope and the selectivity
of the stop bands.

\begin{proposition}[Filter properties]
\label{prop:filter}
For every $\alpha>0$ and every finite $\Omega$ with $\sigma_{m}>0$:
(i) $0\le H_{\alpha,\Omega}(\omega)\le 1$;
(ii) $H_{\alpha,\Omega}(-\omega)=H_{\alpha,\Omega}(\omega)$, hence
$\mathcal{S}_{\alpha,\Omega}$ maps real images to real images;
(iii) $\mathcal{S}_{\alpha,\Omega}$ is self-adjoint and nonexpansive,
$\|\mathcal{S}_{\alpha,\Omega}x\|_{2}\le\|x\|_{2}$;
(iv) $K=0$ recovers the pure fractional-order envelope
$G_{\alpha}(\omega)=\exp(-\alpha|\omega|^{2}/4\omega_{N}^{2})$.
\end{proposition}

\begin{proof}
The envelope lies in $(0,1]$ and each bracket in \eqref{eq:filter} lies in
$[0,1)$, so their product lies in $[0,1]$, which gives (i). Statement (ii)
follows because $|\omega|$ is even and the notch factors are applied in
conjugate pairs. For (iii), $\mathcal{S}_{\alpha,\Omega}$ is diagonal in the
Fourier basis with real eigenvalues $H_{\alpha,\Omega}(\omega)$, hence
self-adjoint, and by (i) its spectral radius is at most one, which by
Parseval gives nonexpansiveness. Setting $K=0$ removes the product, giving
(iv).
\end{proof}

Property (iii) is the reason the method cannot amplify: whatever the data,
the spectral half-step is a contraction, in contrast with state dependent
diffusion coefficients whose gain is unbounded. Section \ref{sec:results}
reports the empirical consequence.

\subsection{Damped Spectral EM Iteration}
\label{sec:iteration}

Let $c=A^{\!\top}\mathbf{1}$ denote the sensitivity image. Starting from
$f^{0}=\mathbf{1}$, each iteration applies a damped multiplicative EM
correction followed by the spectral half-step,
\begin{align}
  f^{k+\frac{1}{2}}
  &= f^{k}\odot\left[\frac{A^{\!\top}\!\left(g\oslash Af^{k}\right)}{c}\right]^{\gamma},
  \label{eq:em}\\[2pt]
  f^{k+1}
  &= \max\!\left\{(1-\lambda)\,f^{k+\frac{1}{2}}
     +\lambda\,\mathcal{S}_{\alpha,\Omega}f^{k+\frac{1}{2}},\;0\right\},
  \label{eq:spectral}
\end{align}
where $\oslash$ is elementwise division, the power in \eqref{eq:em} is
elementwise, $\gamma\in(0,1]$ damps the multiplicative correction and
$\lambda\in[0,1]$ relaxes the spectral step. Setting $\gamma=1$ and
$\lambda=0$ recovers standard MLEM. The exponent form in \eqref{eq:em} keeps
the update multiplicative, so nonnegativity is preserved without projection,
and it is the form for which the mirror-descent interpretation of EM applies;
a convex combination of $f^{k}$ and the EM image is a different iteration and
is not used here.

Throughout we use $\gamma=0.9$ and $\lambda=0.95$, both frozen before any
evaluation on test data.

\subsection{Order Selection by the Discrepancy Principle}
\label{sec:alpha}

The fractional order is selected from the data through the Poisson deviance
between the measured counts and the counts predicted by a candidate
reconstruction. For a reconstruction $f$ with predicted means
$\mu_{i}(f)=(Af)_{i}/\kappa$, the mean deviance per ray is
\begin{equation}
  D(f)=\frac{2}{M}\sum_{i=1}^{M}
  \left[\,y_{i}\log\frac{y_{i}}{\mu_{i}(f)}-\bigl(y_{i}-\mu_{i}(f)\bigr)\right],
  \label{eq:deviance}
\end{equation}
with the convention $y\log y=0$ at $y=0$. Under the model \eqref{eq:poisson},
$D(f^{\star})\approx 1$ for a reconstruction that explains the data to within
the noise, which is the Poisson analogue of Morozov's discrepancy principle.
We therefore select the largest order whose converged fit remains
noise consistent,
\begin{equation}
  \alpha^{\star}=\max\left\{\alpha\in\mathcal{L}:\;D\!\left(f_{\alpha}\right)\le\tau\right\},
  \qquad \tau=1,
  \label{eq:morozov}
\end{equation}
where $f_{\alpha}$ is the reconstruction obtained with order $\alpha$ and no
notches, and the candidate set $\mathcal{L}$ is the geometric ladder \eqref{eq:ladder}
\begin{equation}
  \mathcal{L}=\left\{\alpha_{0}2^{\,k}\right\}_{k=0}^{K_{\max}},
  \qquad \alpha_{0}=0.2,\quad K_{\max}=7 .
  \label{eq:ladder}
\end{equation}
The ladder is evaluated in ascending order and the search stops at the first
violation of $D\le\tau$, which is exact because $D(f_{\alpha})$ is monotone
increasing in $\alpha$: stronger smoothing can only widen the gap between the
model and the data. A bounded grid must not be used here. The order required
at a given operating point grows as the dose falls, and a ceiling silently
turns the selection into a fixed choice; on our data the optimum moves from
$\alpha=1.0$ at $300$ counts per ray to $\alpha=3.2$ at $75$ counts per ray
at the same resolution and sampling. Selection is unaffected by resolution
itself, consistent with the normalisation in \eqref{eq:filter}.

\subsection{Robust Notch Detection}
\label{sec:notch}

Narrowband directional artefacts appear in the reconstruction as isolated
peaks of the magnitude spectrum. Detecting them on clinical anatomy is not a
thresholding exercise: anatomical spectra are strongly anisotropic, so an
isotropic background model produces false positives on every slice. We use a
four stage rule.

\emph{Sectoral background.} The magnitude spectrum
$P=|\mathcal{F}\{f\}|$ is partitioned into cells indexed by radius bin and
angular sector ($S=16$ sectors), and the background $B$ is the median of $P$
within each cell, with a per radius median as fallback for sparsely
populated cells. The detection statistic is the ratio $R=P\oslash B$.

\emph{Robust standardisation.} Over the annulus
$\rho_{0}<\|\omega\|<\rho_{\max}$ we compute the median $m_{0}$ and the
median absolute deviation of $R$, and standardise
\begin{equation}
  z(\omega)=\frac{R(\omega)-m_{0}}{1.4826\,\mathrm{MAD}(R)} .
  \label{eq:zscore}
\end{equation}
Candidates are local maxima of $R$ in a $7\times 7$ neighbourhood whose
standardised score \eqref{eq:zscore} exceeds $\kappa_{z}=10$.

\emph{Acquisition band guard.} Candidates are restricted to
$\rho_{\max}=n_{d}/2-2$, the detector Nyquist radius of the acquisition.
Content beyond this radius cannot have been measured, so a peak there is a
property of the inversion rather than of the object, and notching it removes
nothing that the data contain.

\emph{Spatial coherence gate.} For each surviving candidate, a one sided
band-pass image is formed around $\omega_{m}$ and the fraction of the field
of view where its envelope exceeds $30\%$ of its maximum is measured. A
global stripe or ring covers the field, whereas anatomical texture forms
localised packets; candidates with coverage below $0.35$ are rejected.

Finally, the rule is applied to two iterates of the same reconstruction and
only peaks that persist in both, within three pixels, are retained, which
removes noise induced flukes. The half width of $R$ at half maximum sets
$\sigma_{m}$, capped at $0.05\pi$.

\subsection{Algorithm}
\label{sec:algorithm}

Algorithm \ref{alg:afonf} states the complete method. It has no run time
parameter beyond the frozen constants $\gamma$, $\lambda$, $\tau$ and the
detection thresholds: the order and the notch set are both determined by the
data at hand.

\begin{algorithm}[t]
\caption{A-FONF reconstruction}
\label{alg:afonf}
\begin{algorithmic}[1]
\REQUIRE counts $y$, calibration $\kappa$, system matrix $A$, detector width
$n_{d}$; constants $\gamma=0.9$, $\lambda=0.95$, $\tau=1$, ladder
$\mathcal{L}$, iteration budgets $T_{\mathrm{sel}}$, $T$
\STATE $g\leftarrow\kappa y$, \; $c\leftarrow A^{\!\top}\mathbf{1}$
\STATE \textbf{Stage 1 (order):}
\FOR{$\alpha\in\mathcal{L}$ in ascending order}
  \STATE $f_{\alpha}\leftarrow$ iterate \eqref{eq:em}, \eqref{eq:spectral}
         for $T_{\mathrm{sel}}$ steps with $\gamma=1$, $\Omega=\emptyset$
  \STATE $D_{\alpha}\leftarrow D(f_{\alpha})$ from \eqref{eq:deviance}
  \IF{$D_{\alpha}>\tau$} \STATE \textbf{break} \ENDIF
\ENDFOR
\STATE $\alpha^{\star}\leftarrow\max\{\alpha:D_{\alpha}\le\tau\}$
\STATE \textbf{Stage 2 (notches):}
\STATE $f'\leftarrow$ $20$ steps of \eqref{eq:em}, \eqref{eq:spectral} with
       $\alpha^{\star}$, warm started from $f_{\alpha^{\star}}$
\STATE $\Omega\leftarrow$ peaks detected in both $f_{\alpha^{\star}}$ and
       $f'$ (Section \ref{sec:notch})
\STATE \textbf{Stage 3 (reconstruction):}
\STATE $f\leftarrow$ $T$ steps of \eqref{eq:em}, \eqref{eq:spectral} with
       $(\alpha^{\star},\Omega)$
\RETURN $f$, $\alpha^{\star}$, $\Omega$
\end{algorithmic}
\end{algorithm}

The selection stage dominates the cost, since it reconstructs once per
candidate order; the early stop in \eqref{eq:morozov} keeps this to the
number of ladder entries actually visited, typically four to six. The cost is
incurred once per acquisition and is reported with the run times in
Section \ref{sec:results}.

\subsection{Geometry and Implementation}
\label{sec:implementation}

All operators are used through a single interface, so the same algorithm runs
unchanged in parallel beam and fan beam geometry. Two backends are provided.
A sparse matrix backend assembles $A$ once with bilinear ray weights and
supplies $A^{\!\top}$ as the exact adjoint, which is required for the
monotonicity argument and is verified numerically by
$\langle Ax,u\rangle=\langle x,A^{\!\top}u\rangle$ to machine precision. A
matrix free backend wraps a GPU projector for the clinical scale geometry,
where an explicit matrix is impractical; the unmatched forward and back
projector pair of that backend is used only for the large scale study, and
the two backends are shown to agree by construction of the same iteration.
Fan beam acquisitions use a flat detector with isocentre equispaced bins, so
that the parallel beam matrix is recovered bin for bin as the source distance
grows.

Reference implementations of all baselines, the evaluation code, the fixed
patient splits and the automated verification suite are released with the
paper.
